\documentclass[a4paper,num-refs]{oup-contemporary}
\journal{gigascience}
% Publication metadata remains provisional; no submission venue selected.
\jname{TODO: journal}
\jlogo{}
\title{Marin DNA: TODO title}
\author[1]{TODO: authors}
\affil[1]{TODO: affiliations}
\runningauthor{TODO: authors}
\runningtitle{Marin DNA}
\papercat{TODO: article type}
\jvolume{TODO}
\jnumber{TODO}
\jyear{TODO}
\newcommand{\todo}{\par\noindent\textit{TODO}\par}
\newcommand{\figslot}[2]{%
  \begin{figure}[htbp]
  \centering\fbox{\parbox[c][25mm][c]{0.85\linewidth}{\centering TODO: #1}}
  \caption{#1. TODO}\label{fig:#2}
  \end{figure}}
\newcommand{\tableslot}[2]{%
  \begin{table}[htbp]
  \caption{#1. TODO}\label{tab:#2}
  \centering\begin{tabular}{ll}\toprule
  TODO & TODO\\\midrule
  TODO & TODO\\\bottomrule
  \end{tabular}\end{table}}
\begin{document}
\begin{frontmatter}
\maketitle
\begin{abstract}
\todo
\end{abstract}
\begin{keywords}
TODO
\end{keywords}
\end{frontmatter}
\section{Background}
Genomic language models (gLMs) learn statistical relationships among nucleotides
by predicting DNA sequence from its surrounding context. Training requires
sequences rather than experimentally measured labels, allowing these models
to use reference genomes from multiple species. The resulting sequence
probabilities and representations can support functional constraint prediction,
sequence generation, and adaptation to downstream tasks~\citep{benegas2025review}.
Examples include DNABERT and Nucleotide Transformer for adapting pretrained
representations to genomic prediction tasks~\citep{ji2021dnabert,dallatorre2025nt},
and regLM for regulatory sequence design using autoregressive models together
with supervised predictors~\citep{lal2024reglm}.
Variant effect prediction (VEP) is a particularly useful application: a model
that captures constraints imposed by natural selection may help identify
changes that disrupt biological function. Related approaches to protein
sequence variation have yielded disease-variant predictions without training
on clinical labels~\citep{frazer2021eve}. Accurate, scalable predictions are
needed to interpret rare variants and prioritize candidate causal variants
at associated loci, including changes outside protein-coding
regions~\citep{benegas2023dna,zhou2015deepsea}.
Fine-mapping methods formalize the challenge of resolving candidate causal
variants within associated regions~\citep{benner2016finemap}.

Existing approaches obtain information about variant effects from different
sources. Models trained on functional-genomics measurements predict changes
in molecular phenotypes: DeepSEA predicts chromatin features~\citep{zhou2015deepsea},
Enformer integrates long-range sequence information for expression
prediction~\citep{avsec2021enformer}, and Borzoi predicts RNA-seq
coverage~\citep{linder2025borzoi}. Conservation scores such as
phastCons~\citep{siepel2005phastcons} and phyloP~\citep{pollard2010phylop}
use evolutionary comparisons at homologous sites; GPN-MSA learns nucleotide
predictions from multispecies alignments~\citep{benegas2025gpnmsa}. Both provide valuable
evidence, but their application depends on the availability of experimental
data or suitable whole-genome alignments. Models trained on unaligned DNA
sequences provide another route to learning functional constraint. GPN, for
example, showed that training on the reference genomes of Arabidopsis and
related species can yield informative genome-wide variant scores without
functional-genomics supervision~\citep{benegas2023dna}. Such approaches are
especially attractive for organisms with limited experimental resources,
although success in one species does not establish comparable performance
across different genomes or evolutionary distances.

The composition of the training data is central to this approach. Large
eukaryotic genomes contain abundant repetitive and weakly constrained
sequence, so sampling uniformly need not provide the most useful examples
for learning functional elements~\citep{benegas2025review}.
GPN demonstrated benefits from multispecies
training and from reducing the influence of repetitive sequence on the
training objective~\citep{benegas2023dna}. More generally, the choice of species,
regions, and sequence weights determines which evolutionary and functional
signals a model encounters~\citep{benegas2025review}. Increasing the number
of training bases therefore need not be equivalent to increasing the amount
of information relevant to variant interpretation. Explicit data curation
and sampling offer ways to concentrate training on the capabilities of
interest, but introduce choices that must be evaluated empirically.

TraitGym provides evidence that these choices matter for human regulatory
variants~\citep{benegas2025traitgym}. Its comparisons show substantial gains
with model scale for Evo 2, alongside continuing difficulties with enhancer
variants. Its GPN-Promoter model, trained specifically on animal promoter
sequences, was competitive with Evo 2 40B on the proximal non-exonic
Mendelian subset~\citep{benegas2025traitgym}. These observations motivate a broader question: can one
model retain the complementary capabilities of region specialists through
an appropriate training mixture? Simply pooling datasets gives the largest
region classes the greatest influence. Controlling their relative exposure
instead makes it possible to test how coding, promoter-proximal, and distal
sequence contribute to a shared model, and whether introducing additional
region types improves their corresponding downstream tasks.

Architecture and optimization complicate this comparison. HyenaDNA uses
implicit convolutions to extend genomic context at single-nucleotide
resolution~\citep{nguyen2023hyenadna}. Evo 2 combines a
specialized architecture with extensive data curation and large-scale training,
and provides an important reference for alignment-free variant effect
prediction~\citep{brixi2026evo2}. Here, we ask how far a standard causal,
decoder-only Transformer can progress by improving the data and training
recipe. We use Marin, an open platform for foundation-model development,
to organize data preparation, training, and evaluation~\citep{marin}.
Keeping the architecture within a widely used model family also allows
training and inference to benefit from established software implementations.
Our aim is to assess this practical recipe, rather than to claim that causal
language modeling is a new approach to genomic sequence.

A useful comparison across model sizes also requires an optimization recipe
that remains effective as scale changes. Language-model scaling studies
relate sequence loss to model size, data, and compute~\citep{kaplan2020scaling},
while hyperparameter-transfer methods seek to reduce the cost of tuning
larger networks using small reference models~\citep{yang2021transfer}.
Differences in training quality can obscure the relationship between model capacity, sequence loss,
and downstream performance. We therefore calibrate hyperparameters on small
reference models and test their transfer before studying parameter scaling.
We focus initially on short sequences centered on individual functional
elements, using a 255-base context to make these experiments tractable.
This design supports controlled studies of local sequence constraints; it
does not address the full range of long-distance regulatory interactions.
Together, mixture control and hyperparameter transfer provide a basis for
separating questions about data exposure from questions about model scale.

We evaluate the resulting models using human Mendelian variants and
experimental saturation genome-editing measurements. ClinVar aggregates
clinical interpretations of human variation~\citep{landrum2020clinvar},
while saturation genome editing measures variant effects directly in a
cellular assay, as demonstrated for BRCA1~\citep{findlay2018sge}.
MaveDB provides access to multiplexed variant-effect measurements~\citep{esposito2019mavedb}.
These resources support complementary evaluations: clinical interpretation
and measured cellular function are related but distinct outcomes.
Human data provide
well-characterized tests of biological relevance, even though a longer-term
motivation is to develop useful models for less studied organisms.
We use both sequence log-likelihood ratios and linear probes trained on
frozen representations. Likelihood-based scoring was used in
GPN~\citep{benegas2023dna}; TraitGym compared likelihood and embedding-based
scores on human regulatory variants~\citep{benegas2025traitgym}.
These readouts ask different questions: whether
sequence probabilities rank deleterious changes, and whether the learned
representations contain information useful for that ranking. Neither
should be assumed to summarize all capabilities of a pretrained model.
Our evaluation concerns functional constraint and pathogenicity-related
ranking, rather than direct prediction of changes in gene expression.
Transfer to other species requires separate evaluation.

In this article, we present Marin DNA and examine data mixtures,
hyperparameter transfer, and parameter scaling within this standardized
training framework. We test how changes in sequence loss relate to VEP
under both readouts, and investigate mixtures that add ncRNA and enhancer
sequence to annotation-derived training regions. We show that the resulting
1B-parameter model is competitive with Evo 2 40B on our Mendelian benchmark,
with substantially lower training compute and faster variant scoring at
the models' deployed context lengths. These comparisons characterize a
specific evaluation setting, rather than establish superiority across all
genomic tasks. The experiments assess the extent to which a compact,
standard architecture can combine useful performance with practical training
and inference costs.

\section{Results}
\subsection{Early mixture experiments}
\todo
\figslot{Mixture comparisons}{mixtures}
\subsection{Hyperparameter transfer}
\todo
\figslot{Hyperparameter transfer}{transfer}
\subsection{Parameter scaling}
\todo
\figslot{Loss scaling}{scaling}
\subsection{Downstream performance}
\todo
\figslot{Downstream performance across model scale}{downstream}
\subsection{Later mixture experiments}
\todo
\figslot{Training mixture histories and evaluation}{later-mixtures}
\subsection{Benchmark comparisons and efficiency}
\todo
\figslot{Benchmark comparisons and throughput}{benchmarks}
\section{Discussion}
\todo
\section{Methods}
\subsection{Dataset construction}
\todo
\figslot{Dataset provenance}{datasets}
\subsection{Architecture and training}
\todo
\tableslot{Model configurations}{models}
\subsection{Hyperparameter transfer}
\todo
\subsection{Scaling experiments}
\todo
\subsection{Evaluation benchmarks}
\todo
\figslot{Evaluation design}{evaluation}
\tableslot{Benchmark summaries}{benchmarks}
\subsection{Zero-shot scoring and linear probes}
\todo
\subsection{Efficiency measurements}
\todo
\section{Availability of source code and requirements}
\todo
\section{Data availability}
\todo
\section{Funding}
\todo
\section{Competing interests}
\todo
\section{Acknowledgements}
\todo
\bibliography{paper}
\section*{Supplementary material}
\todo
\end{document}
